\chapter{Useful tools}

\section{Distributions and Hölder regularity}

As we saw in our heuristic regularity power counting, we will have to work with \emph{functions} of negative regularity, i.e. distributions.
Consider the space $\mathscr{D}(\mathbb{R}^d) \coloneq C^{\infty}_c(\mathbb{R}^d)$ of smooth functions with compact support on $\R^d.$
The space of distribution $\mathscr{D}'(\mathbb{R}^d)$ is the dual of $\mathscr{D}(\mathbb{R}^d)$ : a distribution $T$ on $\R^d$ is a linear form $\mathscr{D}(\mathbb{R}^d) \longrightarrow \R$, 
such that for all compact $K \subset \R^d$, there exists $r_k \in \N$ with 
$$\abs{T(\phi)}\lesssim \vert \vert \phi \vert \vert_{C^r} \coloneq \max_{\abs{k}\leq r} \vert \vert D^k \phi \vert \vert_{\infty} \quad \forall \phi \in \mathscr{D} $$

We saw that distributions generalise function because every locally integrable function $f$ defines a distribution by integration :
$$\forall \phi \in \mathscr{D},\quad f(\phi) \coloneq \int_{\R^d} f(x)\phi(x) \dd x .$$

Of course $\mathscr{D}'$ is a vector space. We can also define derivatives of arbitrary order $k = (k_1, \dots, k_d)$ of a distribution $T$ by setting :
$$\forall \phi \in \mathscr{D},\quad D^k T(\phi) \coloneq (-1)^{\abs{k}}T(D^k \phi).$$
when $T$ is a differential function, this corresponds to integration by part.

For $\phi \in C^\infty$ and $T \in \mathscr{D}'$, we can define the product $\phi \cdot T$ by :
$$\forall \psi \in \mathscr{D},\quad (\phi \cdot T) (\psi) \coloneq T(\phi \psi) \quad ( \text{indeed, }\phi \psi \in \mathscr{D}).$$
But in general, there is no canonical way of multiplying two distriubtions. This is really the heart of the problem of singular SPDEs.
To solve this, we will need to formalise a notion of regularity for distributions.

\begin{remark}
    Let $r\in \N$, we denote by $\mathcal{B}^r$ the set of test-function $\phi\in \D$ supported in the unit ball $B(0,1)$ and such that $\norm{\phi}_{C^r}\leq 1$.
\end{remark}

\begin{definition}[Scaling]
    For $\phi \in \D, \lambda \in \R, x \in \R^d$, we set 
    $$\forall y \in \R^d, \quad \phi^\lambda_x(y) \coloneq \lambda^{-d} \phi \left(\frac{y-x}{\lambda}\right)$$
\end{definition}

\begin{remark}
    For $\phi,\psi \in \D$, and $x \in \R^d,\lambda>0$, we have $(\phi*\psi)^\lambda_x = \phi^\lambda * \psi^\lambda_x$.
    
    Indeed,
    \begin{align*}
        \forall y \in \R^d, (\phi*\psi)^\lambda_x (y) &= \lambda^{-d} \int\dd z \phi(z) \psi(\lambda^{-1}(y  - x) - z) \\
        &= \int \dd z \lambda^{-d} \phi(\lambda^{—1}z) \lambda^{-d} \psi(\lambda^{-1}(y - x - z)) \quad \text{(change of variable }z\leftarrow \lambda z) \\
        &= \int \dd z \phi^\lambda(z) \psi^\lambda_x(y - z).
    \end{align*}
\end{remark}

Another operation on distribution that will be very useful is the convolution by a test-function.

\begin{definition}
    Let $T\in \D'$ and $\rho \in \D$. Define the function $T*\rho$ by :
    \begin{equation*}
        \forall x \in \R^d, \quad (T*\rho)(x) = T(\rho(x - \cdot)).
    \end{equation*}    
\end{definition}

\begin{proposition}
    With the same notation as in the definition, $T*\rho$ is smooth and for all $k \in \N^d$, 
    \begin{equation*}
        D^k(T*\rho) = (D^kT) * \rho = T*(D^k\rho)
    \end{equation*}
\end{proposition}

\begin{proof}
    The second equality holds by definiton of the derivative of a distribution. So to prove the proposition, 
    it suffices to show that $T*\rho$ is differentiable and that for $i \in \{1,\dots,d\}, \quad \partial^{i}(T*\rho) = T*(\partial^{i}).$
    \begin{equation*}
        \forall x,y \in \R^d, \quad T*\rho(y) - T*\rho(x) - \sum{i=1}{d} (y_i - x_i) T*(\partial^{i}\rho)(x) = T(\eta_{x,y}),
    \end{equation*}
    where $\eta_{x,y} = \rho(y - \cdot) - \rho(x- \cdot) - \sum_{i=1}^{d}(y_i - x_i)\partial^{i}\rho(x- \cdot)$.
    Fix $x\in \R^d$ and $R>0$. There exists a compact $K$ such that for all $y\in B(x,R)$, the support of $\eta_{x,y}$ is included in $K$.
    So by definition of a distribution, there exists $r\in \N$ such that :
    $$\forall y \in B(x,R), \quad \abs{T(\eta_{x,y})}\lesssim \norm{\eta_{x,y}}.$$
    Now by Taylor's theorem, $\norm{\eta_{x,y}} = o(x-y)$, so $T*\rho$ is differentiable and $\partial^{i}(T*\rho) = T*(\partial^{i}\rho)$ for $i\in \{1,\dots,d\}.$
\end{proof}

\begin{lemma}
    Let $\phi \in D$ be a test-function. $T*\rho$ viewed as a distribution ($C^\infty \subset L^\infty_{\text{loc}}$),
    $$(T*\rho) (\phi) = T(\tilde{\rho}*\phi),$$
    where $\tilde{\rho} = \rho(-\cdot)$.
\end{lemma}

\begin{proof}
    Fix $a<b\in \R$ such that $\operatorname{Supp}(\phi)\subset [a,b]^d.$
    Let $\delta_n \coloneq \frac{b-a}{2^n}$ and $\mathcal{P}_n \coloneq \{a+i\delta_n \vert i = 0, \dots, 2^n - 1\}.$
    \begin{align*}
        (T*\rho) (\phi) &= \int (T*\rho)(x) \phi(x) \dd x \\
        &= \int T(\rho(x- \cdot)) \phi(x) \dd x \\
        &= \int_{[a,b]^d} T(\rho(x- \cdot)\phi(x)) \dd x \\
        &= \lim_{n \to +\infty} T(\underbrace{\sum_{x \in \mathcal{P}_n^d}\delta_n^d \rho(x- \cdot)\phi(x)}_{I_n(\cdot)}) \quad \text{(Riemann sum)}\\
        &= T(\lim_{n \to +\infty} \sum_{x \in \mathcal{P}_n^d}\delta_n \rho(x- \cdot)\phi(x)) \\
        &= T(\int\rho(x- \cdot)\phi(x) \dd x)\\
        &= T(\tilde{\rho}*\phi).
    \end{align*}
    The exchange $\lim_n T(I_n) = T(\lim_n I_n)$ is justified by the facts that :
    \begin{itemize}
        \item because $[a,b]^d$ and $\operatorname{Supp}(\eta)$ are compact, there exists a compact $K$ such that for all $n\in \N$, $\operatorname{Supp}(\sum_{x \in \mathcal{P}_n^d}\delta_n^d \rho(x- \cdot)\phi(x))\subset K$,
        \item by definition of a distribution, there exists $r\in \N$ such that T is continuous on $C^\infty(K)$ equiped with $\norm{\cdot}_{C^r}$,
        \item finally, $\sum_{x \in \mathcal{P}_n^d}\delta_n^d \rho(x- \cdot)\phi(x) \xrightarrow[n\to + \infty]{\norm{\cdot}_{C^r}} \int\rho(x- \cdot)\phi(x) \dd x $ by the Riemann sum property on the successive derivatives.
    \end{itemize}
\end{proof}

\begin{theorem}[Mollification]
    Let $T \in \D'$ be a distribution, and $\rho \in \D$ be any test-function such that $\int \rho = 1$.
    Then, for any test-function $\phi \in \D$, and all $r\in \N$,
    \begin{equation*}
        \norm{\phi - \rho^\lambda * \phi}_{C^r} \xrightarrow[\lambda \to 0]{} 0,
    \end{equation*}
    and consequently,
    \begin{equation*}
        T*\rho^\lambda(\phi) \xrightarrow[\lambda \to 0]{} T(\phi).
    \end{equation*}
\end{theorem}
First, let $r\in \N$ and $k\in \N^d$ such that $\abs{k}\leq r$.
Differentiating under the integral sign, and knowing $\int \rho = 1$, we have :
\begin{equation*}
    \partial^k(\phi - \rho^\lambda * \phi)(x) = \int \rho^\lambda(z) (\partial^k \phi(x)- \partial^k \phi(z - x)) \dd z.
\end{equation*}
So that by Taylor's theorem,

\begin{align*}
    \norm{\phi - \rho^\lambda * \phi}_{C^r} &\leq \norm{\phi}_{C^{r+1}} \int \rho^\lambda(z) \abs{z} \dd z \\
    &\leq \lambda \norm{\phi}_{C^{r+1}} \int \rho(u) \abs{u} \dd u.
\end{align*}

Now for the second statement, we proved that $T*\rho^\lambda(\phi) = T(\tilde{\rho}^\lambda*\phi)$ and $\tilde{\rho}^\lambda*\phi \xrightarrow[\lambda \to 0]{\norm{\cdot}_{C^r}} \phi$.
Additionaly, we can find a compact $K$ such that for all $\lambda\in (0,1],\quad \operatorname{Supp}(\tilde{\rho}^\lambda*\phi)\subset K$, 
So as before, the continuity property of $T$ applies.


\begin{definition}[$\alpha$-Hölder functions]
    Let $\alpha \in (0,1)$. We say that a function $f$ is $\alpha$-Hölder if 
    $$\forall x,y \in \R^d, \abs{f(x)- f(y)} \lesssim \abs{x-y}^\alpha.$$
    We say that $f$ is locally $\alpha$-Hölder if the constant in the estimate depends on the compact in which we take $x$ and $y$.
    For $\alpha \in \R_+ \setminus \N$, we say that $f$ is $\alpha$-Hölder if f is $\lfloor \alpha \rfloor$-times differentiable, and
    $$\norm{f}_\alpha \coloneq \norm{f}_{\infty} + [f]_\gamma < \infty $$
    where 
    $$[f]_\gamma \coloneq \sup_{x\in \R^d} \sup_{h\in \R^d} \frac{\abs{f(x+h) - \sum_{\abs{k}<\gamma} \frac{h^k}{k!}\partial^kf(x)}}{\abs{h}^\gamma}.$$
\end{definition}

We can also define Hölder spaces for negative $\alpha$.

First, given a positive integer $r$, note $B_r$ the set of smooth functions $\phi$ supported in the unit ball, with :
$$\norm{\phi}_{C^r}\coloneq \sup_{\abs{k}\leq r} \norm{D^k \phi}_\infty \leq 1.$$
\begin{definition}[Negative Hölder spaces]
    Let $\alpha<0$. $T$ is $\alpha$-Hölder if 
    $$\forall \phi \in B_r, \quad \abs{T(\phi^\lambda_x)} \lesssim \abs{\lambda}^\alpha \quad \text{uniformly in } x\in \R^d,\lambda \in (0,1)$$
    where $r>\lceil - \alpha \rceil$.
    We can also talk about local $\alpha$-Hölder distribution, when the estimate is uniform in $x \in K$ for all $K$ compact.
    
    It will be useful to further generalise this notion by relaxing the uniform control in $x$ and $\lambda$, asking for 
    $$\norm{\norm{ \sup_{\phi \in B_r}\frac{T(\phi^\lambda_x)}{\abs{\lambda}^\alpha}}_{L^p(x\in \R^d)}}_{L^q(\lambda\in(0,1))} < \infty$$
    This defines the Besov space $\mathcal{B}^\alpha_{p,q}$.
\end{definition}


\begin{remark}
    \begin{itemize}
        \item We have $\mathcal{B}^\alpha_{p,q_1} \hookrightarrow \mathcal{B}^\alpha_{p,q_2}$ for $q_1 \leq q_2.$
        \item We have $\mathcal{B}^\alpha_{p,q} \hookrightarrow \mathcal{B}^\beta_{p,q}$ for $\beta \leq \alpha$ on compact domains.
    \end{itemize}
\end{remark}


A very useful lemma allows us to check the finiteness of this scaling norm only for one compactly supported test function $\phi$ with $\int \phi \neq 0.$ and for $\lambda$ on the dyadic scale.
\begin{lemma}
    \label{lem:technique1}
    Let $\alpha<0$. Then $T \in \mathcal{B}^\alpha_{p,q}$ if and only if there exists a compactly supported test-function $\phi$ with $\int \phi \neq 0$ and
    $$\norm{\norm{\frac{T(\phi^{2^{-n}}_x)}{2^{-n\alpha}}}_{L^p(x\in \R^d)}}_{l^q(n)} < \infty.$$
\end{lemma}
 We present the proof because it displays ideas that will reappear later, in reconstruction arguments for instance.
\begin{proof}
    The direct implication is clear. 
    Fix $T \in D'$ and $\phi$ such that the statement holds, and let's prove that $T \in \mathcal{B}^\alpha_{p,q}$.
    Fix $r\in \N$ such that $r>-\alpha$, $\lambda_0, \dots,\lambda_{r-1} \in \R^*_+$ distincts, $c_0,\dots,c_{r-1} \in \R$, and
    \begin{equation*}
        \hat{\phi} \coloneq \frac{1}{\int \phi} \sum_{i=0}^{r-1} c_i \phi^{\lambda_i}
    \end{equation*}
    We can tune the constants $c_i$ so that 
    \begin{itemize}
        \item $\int \hat{\phi} = 1$,
        \item $\forall k, 1\leq \abs{k}\leq r-1, \quad \int x^k \hat{\phi}(x) \dd x = 0$, 
    \end{itemize}
    Indeed, 
    \begin{align*}
        \forall k \in \N^d, \abs{k}\leq r-1, \int x^k \hat{\phi}(x) \dd x = \delta_{0,k} &\iff \frac{1}{\int \phi} \sum_{i=0}^{r-1} c_i \int x^k \phi^{\lambda_i} (x) \dd x = \delta_{0,k} \\
        & \iff \frac{1}{\int \phi} \sum_{i=0}^{r-1} c_i \lambda_i^{\abs{k}} \int \phi = \delta_{0,k}  \quad \text{(change of variable)}\\
        & \iff \sum_{i=0}^{r-1} c_i \lambda_i^{\abs{k}} = \delta_{0,k}
    \end{align*}

    Because the $\lambda_i$ are distincts, this is an inversible Vandermonde system.
    Additionaly, by choosing small $\lambda_i$s, we see that we can shrink the support of of $\hat{phi}$ so that it is included in $B(0,\frac{1}{2})$.
    
    
    
    Let $\lambda_i = 2^{m_i}$ for $m_i\in \N$ large enough.
    
    \begin{align*}
        \norm{\norm{\frac{T(\hat{\phi}^{2^{-n}}_x)}{2^{-n\alpha}}}_{L^p(x\in \R^d)}}_{l^q(n)} &\lesssim \sum_{i=0}^{r-1} \norm{\norm{\frac{T(\phi^{2^{-n-m_i}}_x)}{2^{-n\alpha}}}_{L^p(x\in \R^d)}}_{l^q(n)} \quad \text{(triangle inequality)}\\
        &\lesssim \sum_{i=0}^{r-1} \norm{\norm{\frac{T(\phi^{2^{-n-m_i}}_x)}{2^{-(n+m_i)\alpha}}}_{L^p(x\in \R^d)}}_{l^q(n)} \\
        &<\infty
    \end{align*}

    Let $\psi \in D'$ be any test function,and $\rho$ a mollifier with $\int \rho = 1$.
    We know that $\int T(\rho^{2^{-m}}_z) \psi (z)  \dd z  = T(\rho^{2^{-m}}*\psi) \xrightarrow[m \to \infty]{} T(\psi)$.
    So we can write the telescopic sum :
    \begin{equation*}
        T(\psi^{2^{-n}}_x) = \int T(\rho^{2^{-n}}_z) \psi^{2^{-n}}_x (z) \dd z + \sum_{m\geq n} \int T((\rho^{2^{-m-1}} - \rho^{2^{-m}})_z) \psi^{2^{-n}}_x (z) \dd z \\
    \end{equation*}
    Now we use $\rho \coloneq \hat{\phi}^2*\hat{\phi}$, motivated by the fact that we can write the increment in convolution form : $\rho^{2^{-m-1}} - \rho^{2^{-m}} = (\rho^{\frac{1}{2}} - \rho)^{2^{-m}} = \hat{\phi}^{2^{-m}}*\check{\phi}^{2^{-m}}$ where $\check{\phi}= \hat{\phi}^{\frac{1}{2}} - \hat{\phi}^2$.
    Plugging in the equation :
    \begin{align*}
        T(\psi^{2^{-n}}_x) &= \int T(\hat{\phi}^{2^{-n}}*\hat{\phi}^{2^{-n+1}}_z) \psi^{2^{-n}}_x (z) \dd z + \sum_{m\geq n} \int T(\hat{\phi}^{2^{-m}}*\check{\phi}^{2^{-m}}_z) \psi^{2^{-n}}_x (z) \dd z \\
        &= \int T(\hat{\phi}^{2^{-n}}_z) (\hat{\phi}^{2^{-n+1}}*\psi^{2^{-n}}) (z - x) \dd z + \sum_{m\geq n} \int T(\hat{\phi}^{2^{-m}}_z) (\check{\phi}^{2^{-m}}*\psi^{2^{-n}}) (z - x) \dd z \\
        &= \int_{B(0,2^{-n+1})} T(\hat{\phi}^{2^{-n}}_{y+x}) (\hat{\phi}^{2^{-n+1}}*\psi^{2^{-n}}) (y) \dd y + \sum_{m\geq n} \int_{B(0,2^{-n+1})} T(\hat{\phi}^{2^{-m}}_{y+x}) (\check{\phi}^{2^{-m}}*\psi^{2^{-n}}) (y) \dd y  \\
    \end{align*}

    The last line is justified by the fact that $y =z - x$, and that $y \mapsto T(\hat{\phi}^{2^{-n}}_{y+x})$ is supported in $B(0,2^{-n+1})$.
    To continue, we need the following lemma :
    \begin{lemma}
    \label{lem:technique2}
    Assume that $\nu \in D$ is a test function that cancels polynomials up to degree $r-1 \in \N$ : $\forall 0\leq k \leq r-1, \int x^k \nu(x) \dd x =0$.
    Then for any test function $\eta \in D$ and $\lambda >0$,
    \begin{equation*}
        \norm{\nu^\lambda *\eta }_\infty \lesssim \norm{\eta}_{C^r} \norm{\nu}_{L^1} \lambda^r
    \end{equation*}
    where the constant depends on the support of $\nu$ and $r$.
    \end{lemma}
    \begin{proof}
        Fix $\eta \in D$ and $\lambda >0$.
        Denote for all $x,y \in \R^d, \quad T_x^{r-1}\eta (y) = \sum_{\abs{k}\leq r-1 } \frac{D^k \eta (x)}{k!} (y-x)^k$ the Taylor polynomial at $x$ up to order $r-1$ of $\eta$.
        
        By Taylor's theorem, $\abs{(\eta - T_x^r\eta)(y)} \leq \frac{\norm{\eta}_{C^r}}{r!} \abs{y-x}^r$.

        \begin{align*}
            \forall x \in \R^d,\quad \abs{\nu^\lambda *\eta (x)} &\leq \abs{\int \nu^\lambda (y) (\eta - T_x^r\eta)(x-y) \dd y + \underbrace{\int \nu^\lambda (y) (T_x^r\eta)(x-y) \dd y}_{=0 \text{ by hypothesis}}} \\
            &\leq \int \abs{\nu^\lambda (y)} \abs{(\eta - T_x^r\eta)(x-y)} \dd y \\
            &\leq \int \abs{\nu^\lambda (y)} \frac{\norm{\eta}_{C^r}}{r!} \abs{y}^r \dd y \\
            &\leq \frac{\norm{\eta}_{C^r}}{r!} \lambda^r \int_K \abs{\nu (z)} \abs{z}^r \dd z \\
            &\leq \frac{1}{r!}(\sup_{z\in K} \abs{z}^r)  \norm{\eta}_{C^r} \norm{\nu}_{L^1} \lambda^r \quad (\text{uniform bound in }x)
        \end{align*}
        where $K$ denotes the (compact) support of $\nu$.
    \end{proof}

    Thanks to this, we have :
    \begin{empheq}[left=\empheqlbrace]{align*}
    \norm{\hat{\phi}^{2^{-n+1}}*\psi^{2^{-n}}}_\infty &\lesssim 2^{nd} \norm{\psi}_\infty \norm{\hat{\phi}}_{L^1} \leq2^{nd} \norm{\psi}_{C^r} \norm{\hat{\phi}}_{L^1}  \\
    \norm{\check{\phi}^{2^{-m}}*\psi^{2^{-n}}}_\infty &\lesssim 2^{nd} 2^{(n-m)r} \norm{\psi}_{C^r} \norm{\check{\phi}}_{L^1}
    \end{empheq} 
    As a consequence, denoting $B_n = B(0,2^{-n})$,
    \begin{equation*}
        \sup_{\psi \in B^r} \abs{T(\psi^{2^{-n}}_x)} \lesssim 2^{nd} \norm{\hat{\phi}}_{L^1} \int_{B_{n-1}} T(\hat{\phi}^{2^{-n}}_{y+x}) \dd y + \sum_{m\geq n} 2^{nd + (n-m)r} \norm{\check{\phi}}_{L^1} \int_{B_{n-1}} T(\hat{\phi}^{2^{-m}}_{y+x}) \dd y.
    \end{equation*}
    Taking $L^p$-norm in $x$, we get by triangle inequality :
    \begin{equation*}
        \norm{\sup_{\psi \in B^r} \abs{T(\psi^{2^{-n}}_x)}}_{L^p(x)} \lesssim 2^{nd} \norm{\hat{\phi}}_{L^1} \int_{B_{n-1}} \norm{T(\hat{\phi}^{2^{-n}}_{y+x})}_{L^p(x)} \dd y + \sum_{m\geq n} 2^{nd + (n-m)r} \norm{\check{\phi}}_{L^1} \int_{B_{n-1}} \norm{T(\hat{\phi}^{2^{-m}}_{y+x})}_{L^p(x)} \dd y.
    \end{equation*}
    Now in these integrals, the integrand is actually constant, we get :
    \begin{equation*}
        \norm{\sup_{\psi \in B^r} \abs{T(\psi^{2^{-n}}_x)}}_{L^p(x)} \lesssim \norm{T(\hat{\phi}^{2^{-n}}_{x})}_{L^p(x)} + \sum_{m\geq n} 2^{(n-m)r} \norm{T(\hat{\phi}^{2^{-m}}_{x})}_{L^p(x)}.
    \end{equation*}
    Then, dividing by $2^{-n\alpha}$ and taking the $l^q$-norm in $n$,

     \begin{equation*}
        \norm{T}_{\mathcal{B}^\alpha_{p,q}} \lesssim \norm{\norm{\frac{T(\hat{\phi}^{2^{-n}}_{x})}{2^{-n\alpha}}_{L^p(x)}}}_{l^q(n)} +  \norm{\sum_{m\geq n} 2^{(n-m)(r+\alpha)} \norm{\frac{T(\hat{\phi}^{2^{-m}}_{x})}{2^{-m\alpha}}_{L^p(x)}}}_{l^q(n)}.
    \end{equation*}
    And finally, because we chose $r>-\alpha$, applying Jensen's inequality, then using Fubini for the sums in $n$ and $m$, we proved that $\norm{T}_{\mathcal{B}^\alpha_{p,q}}<\infty$.





    
\end{proof}


Juste like with Sobolev spaces with Sobolev embeddings, we can trade regularity for integrability via the following theorem.

\begin{theorem}
    Let $q>0$, $0<p_0 \leq p_1 \leq \infty$, and $-\infty < s_1 \leq s_0 < \infty$.
    Then,
    $$s_0 - \frac{n}{p_0} = s_1 - \frac{n}{p_1} \implies \mathcal{B}^{s_0}_{p_0,q} \hookrightarrow \mathcal{B}^{s_1}_{p_1,q}.$$
\end{theorem}

See \cite{Tri83} for a general proof. We are only going to use this in a specific case where the proof will be easier.

We said that we could not multiply two distributions in general, but Hölder regularity gives a simple sufficient condition.


    



\section{Random distributions and noise}

\begin{definition}[Random distributions]
    A random distribution $T$ is a map $(\Omega, \mathcal{F}, \mathbb{P}) \longrightarrow \D'$ such that for all test-function $\phi \in \D$,
    $T(\phi)\coloneq T \circ \text{ev}_\phi : (\Omega, \mathcal{F}, \mathbb{P}) \longrightarrow \R$ is measurable (i.e. a real random variable).
\end{definition}

The forcing term $\xi$ in our equation is a random distribution. A common choice is white noise, let us describe it's law.


\section{Germs and reconstruction}

We return to our problem of defining the product between a distribution and a function. 
Assume that our function $f$ has (positive) regularity $\alpha$, and our distribution $g$ negative regularity $\beta$.
If it is not clear how we could define the product $fg$ globally, we know that around a point $x$, $f$ is well-approximated by its Taylor polynomial at $x$ : $T_x^\alpha f = \sum_{\abs{k}<\alpha}\frac{f^{(k)(x)}}{k!} (\cdot - x)^k$.
Of course as a polynomial, $T_x^\alpha f$ is smooth, so the product $(T_x^\alpha f) g$ is well-defined, and is expected to approximate the formal $fg$ around $x$.
The way of defining $fg$ will be to reconstruct a global distribution from the local approximations $T_x^\alpha f$, for $x \in \R^d$.

\begin{definition}[Germs of distribution]
    A germ of distribution is a family of distribution $(F_x)_{x \in \R^d}$ indexed by our space variable. Additionaly, we require for technical reason that for all test-function $\phi \in \D$, the map $x\mapsto F_x(\phi)$ should be measurable.
\end{definition}

\begin{example}
    $((T_x^\alpha f)g)_{x\in \R^d}$ is a germ of distributions.
\end{example}


\begin{theorem}[Reconstruction]
    \label{thm:germ_reconstruction}
    Let $(F_x)_{x \in \R^d}$ be a germ. Let $\gamma > 0$, and assume that for some $\phi \in \D$ with $\int \phi \neq 0$ and some $\alpha < \gamma$, we have the coherence condition :
    \begin{equation}
        \label{eq:coherence1}
        \norm{F}_{\alpha,\gamma;\phi} \coloneq \sup_{x,y \in \R^d} \sup_{\lambda >0 } \; \frac{\abs{(F_x - F_y)(\phi^\lambda_x)}}{\lambda^\alpha (\lambda + \abs{x-y})^{\gamma - \alpha}} < \infty.
    \end{equation}
    Then, there exists a unique distribution $\mathcal{R}(F) \in \D'$ called the reconstruction of $F$, that is locally well approximated by the germ :
    \begin{equation*}
        \sup_{x\in \R^d} \sup_{\lambda > 0} \sup_{\psi \in \mathcal{B}^r} \frac{\abs{(\mathcal{R}(F)-F_x)(\psi_x^\lambda)}}{\lambda^\gamma} \lesssim \norm{F}_{\alpha,\gamma;\eta},
    \end{equation*}
    with $r\in \N, r>-\alpha$. The implicit constant depending only on $\alpha, \gamma, r,\phi$ and $d$.
\end{theorem}

\begin{remark}
    \begin{itemize}
        \item The coherence condition \eqref{eq:coherence1} can be interpreted as follows : the regularity of the increments $F_x - F_y$ is $\alpha$, and when setting $\abs{x-y}\leq \lambda \rightarrow 0$, the increment tested on $\eta^\lambda_x$ behaves like $\lambda^\gamma$. This quantity goes to $0$ because $\gamma >0$, this is a continuity condition on $x \mapsto F_x$.
        \item This is one version of a reconstruction theorem, in negative Hölder space, but we will present another example in Besov spaces.
        \item If we have also an homogeneity bound on the germ, meaning there exists $\beta \in \R, \; \alpha<\beta\leq\gamma$ such that $\norm{F}_{\text{coh,}\beta}\coloneq \sup_{x\in \R^d} \sup_{\lambda > 0} \sup_{\psi \in \mathcal{B}^r} \frac{\abs{F_x(\psi_x^\lambda)}}{\lambda^\beta} < \infty$, then by triangle inequality, we see that :
        \begin{equation*}
            \sup_{x\in \R^d} \sup_{\lambda > 0} \sup_{\psi \in \mathcal{B}^r} \frac{\abs{\mathcal{R}(F)(\psi_x^\lambda)}}{\lambda^\beta} \leq \norm{F}_{\text{hom,}\beta} + \norm{F}_{\alpha,\gamma;\phi} < \infty.
        \end{equation*}
        Thus the reconstruction is $\beta \wedge \gamma$-Hölder.
    \end{itemize}
\end{remark}

We return to the problem of mutliplying a distribution with a non-smooth function.
\begin{corollary}[Young's multiplication]
    Let $\alpha \in \R^*_+\setminus\N$, $\beta <0$ be such that $\alpha+\beta>0$. 
    Then, there exists a unique bilinear map $C^\alpha \times C^\beta \longrightarrow C^\beta$, that extends the multiplication operator $C^\infty \times \D' \longrightarrow \D'$.
\end{corollary}

\begin{proof}
    Let $f\in C^\alpha$ and $g\in C^\beta$. Recall that we introduced the germ $(F_x \coloneq (T_x^\alpha f)g)_{x\in \R^d}$ and that we aimed to solve the problem of multiplication by applying reconstruction to it.
    We will check that this germ is $(\beta, \alpha+\beta)$-coherent :
    \begin{align*}
        \forall x,y \in R^d, \quad F_x - F_y &= \left( \sum_{\abs{k}<\alpha} \frac{D^kf(x)}{k!} (\cdot - x)^k - \sum_{\abs{k}<\alpha} \frac{D^kf(y)}{k!} (\cdot - y)^k \right) g \\
        &= \left( \sum_{\abs{k}<\alpha} \frac{D^kf(x)}{k!} (\cdot - x)^k - \sum_{\abs{k}<\alpha} \frac{D^kf(y)}{k!} \sum_{\abs{l}\leq \abs{k}} \binom{k}{l} (\cdot -x)^{l}(x-y)^{k-l} \right) g \\
        &= \left( \sum_{\abs{k}<\alpha} \left(\frac{D^kf(x)}{k!} - \sum_{\abs{k}\leq \abs{j} < \alpha} \frac{D^jf(y)}{j!} \binom{j}{k} (x-y)^{j-k} \right) (\cdot - x)^k \right) g \\
        &= \left( \sum_{\abs{k}<\alpha} \frac{1}{k!} \left(D^kf(x) - \sum_{\abs{j} < \alpha - \abs{k}}  \frac{D^{j+k} f(y)}{j!} (x-y)^{j} \right) (\cdot - x)^k \right) g \\
    \end{align*}
    Now f is in $C^\alpha$, so for $k$ such that $\abs{k}<\alpha$, $D^kf$ is in $C^{\alpha - \abs{k}}$, and
    $$\forall x,y \in \R^d, \quad\abs{D^kf(x) - \sum_{\abs{j} < \alpha - \abs{k}}  \frac{D^{j+k} f(y)}{j!} (x-y)^{j}} \leq [f]_{\alpha} \abs{x-y}^{\alpha -\abs{k}}.$$
    As a result, for any test-function $\eta\in \D$, $\lambda>0$ and $x,y\in \R^d$,
    \begin{align*}
        \abs{(F_y - F_x)(\eta^\lambda_x)} &\leq \sum_{\abs{k}<\alpha} \frac{1}{k!} [f]_{\alpha} \abs{x-y}^{\alpha -\abs{k}} \abs{g((\cdot - x)^k \eta^\lambda_x)} \\
        &\leq \sum_{\abs{k}<\alpha} \frac{1}{k!} [f]_{\alpha} \abs{x-y}^{\alpha -\abs{k}} \lambda^{\abs{k}} \abs{g(((\cdot)^k \eta)^\lambda_x)} \\
        &\leq [f]_{\alpha} \norm{g}_\beta \sum_{\abs{k}<\alpha} \frac{1}{k!} \abs{x-y}^{\alpha -\abs{k}} \lambda^{\beta + \abs{k}}  \quad (\text{by definition of }\beta\text{-Hölder norm})\\
        &\lesssim [f]_{\alpha} \norm{g}_\beta \lambda^\beta (\lambda + \abs{x-y})^{\alpha+\beta}.
    \end{align*}
    
    $\alpha+\beta>0$ so we can apply the reconstruction theorem : there exists a unique $fg \coloneq \mathcal{R}(F)\in \D'$, such that
    \begin{equation*}
        \sup_{x\in \R^d} \sup_{\lambda > 0} \sup_{\phi \in \mathcal{B}^r} \frac{\abs{(fg - (T_x^\alpha f) g)(\phi_x^\lambda)}}{\lambda^{\alpha+\beta}} \lesssim \norm{F}_{\alpha,\gamma;\eta}
    \end{equation*}

    Also, 

    \begin{align*}
        \abs{(F_x)(\eta^\lambda_x)} &\leq \sum_{\abs{k}<\alpha} \abs{\frac{D^kf(x)}{k!}} \abs{g((\cdot - x)^k\eta^\lambda_x)}\\
        &\leq \sum_{\abs{k}<\alpha} \abs{\frac{D^kf(x)}{k!}} \lambda^{\abs{k}} \abs{g(((\cdot)^k\eta)^\lambda_x)}\\
        &\leq \norm{g}_\beta \sum_{\abs{k}<\alpha} \abs{\frac{D^kf(x)}{k!}} \lambda^{\abs{k}+\beta} \\
        &\lesssim \norm{g}_\beta \norm{f}_{C^\alpha} \lambda^{\beta} \quad \text{uniformly in }\eta, \; x \text{ and }\lambda.
    \end{align*}
    So by the third item in the remark above, $fg$ is $\beta$-Hölder.
    Additionaly, the map, $(f,g)\mapsto (T_x^\alpha f) g$ is bilinear, and $\mathcal{R}:F_x \mapsto \mathcal{R}(F)$ is linear (by easily seen by uniqueness). So $(f,g)\mapsto fg$ is bilinear.
    It remains to check that it extends the multiplication map. For this, it is enough to check that the product verifies the reconstruction bound when $f \in C^\infty$.
    
    Suppose $f \in C^\infty$.

    \begin{align*}
        \forall x \in \R^d, \lambda>0, \phi \in \mathcal{B}^r, \abs{(fg - (T_x^\alpha f)g)(\phi^\lambda_x)} &= \abs{g((f - (T_x^\alpha f))(\phi^\lambda_x))} \\
        &= \abs{g(((\sum_{\abs{k} = \lfloor \alpha \rfloor + 1}R_k(\cdot)(\cdot - x)^k ))(\phi^\lambda_x))} \\
        &= \lambda^{\lfloor \alpha \rfloor + 1} \abs{g((((\sum_{\abs{k} = \lfloor \alpha \rfloor + 1}R_k(\cdot + x)^{\lambda^{-1}}(\cdot)^k ))(\phi))^\lambda_x)} \\
        &\lesssim \lambda^{\alpha + \beta} \quad (g \in C^\beta)
    \end{align*}
    The second line is justified by Taylor's theorem because $f$ is smooth, $R_k(\cdot)$ is smooth and vanishes at $x$.
\end{proof}

\begin{proof}
    Fix $(F_x)_{x\in \R^d}$ and $\phi$ as the statement of the theorem.

\textit{Step 1 : Uniqueness}

Let $R_1$ and $R_2$ be two distributions such that the reconstruction bound holds :
\begin{equation*}
        \sup_{x\in \R^d} \sup_{\lambda > 0} \sup_{\psi \in \mathcal{B}^r} \frac{\abs{(R_i-F_x)(\psi_x^\lambda)}}{\lambda^\gamma} \lesssim \norm{F}_{\alpha,\gamma;\phi}, \quad i=1,2.
\end{equation*}
Setting $R = R_1 - R_2$, by triangle inequality, we have :
\begin{equation*}
        \sup_{x\in \R^d} \sup_{\lambda > 0} \sup_{\psi \in \mathcal{B}^r} \frac{\abs{R(\psi_x^\lambda)}}{\lambda^\gamma} <\infty .
\end{equation*}
In particular, because $\gamma>0$, for $\psi \in \mathcal{B}^r, \int \psi = 1$ a mollifier, $\sup_{x\in \R^d} \abs{R(\phi_x^\lambda)}\xrightarrow[\lambda \to 0]{} 0.$

Let $\eta \in \D$ be any test-function.
\begin{align*}
    R(\eta) &= \lim_{\lambda \to 0} \int R(\psi^\lambda_x) \eta(x) \dd x \quad \text{(mollification)}\\
    &=0 \quad \text{(by uniform convergence)}
\end{align*}
So we proved that $R=0$.

\textit{Step 2 : Existence}

We want to construct a distribution well-approximated by $F_x$ close to $x$. Let $\rho \in \D$ be a test-function such that $\int \rho =1$, to be specified later.
By analogy with the mollification property, we want to define for all $n\in \N$ and $\eta \in \D$, $R_n (\eta) \coloneq \int F_x * \rho^{2^{-n}}(x)\eta(x) \dd x$ and hope to show that $R_n$ has a limit when $n$ goes to $+\infty$.

The function $R_n : x\mapsto F_x * \rho^{2^{-n}}(x)$ is well-defined and measurable. Indeed, for all $y\in \R^d, \quad x \mapsto F_x*\rho^{2^{-n}}(y) =F_x(\rho^{2^{-n}}(y - \cdot)) $ is measurable by definition of a germ, and for all $x\in \R^d, \quad y\mapsto F_x*\rho^{2^{-n}}(y)$ is continuous by the mollification property.

We want to estimate the difference $\mathcal{R}(F) - F_x$ so it is natural to introduce the mollified version : 
\begin{equation*}
    \forall x,z \in \R^d,\;n\in \N, \quad f_{x,n} (z) = R_n(z) - F_x*\rho^{2^{-n}}(z) = (F_z - F_x)*\rho^{2^{-n}}(z).
\end{equation*}
Fix $x\in \R^d$. $f_{x,n}$ is a measurable function. We want to show that it converges to a distribution when $n$ goes to $+\infty$.

\begin{equation*}
    \forall z\in \R^d, \quad (f_{x,n+1} - f_{x,n})(z) = (F_z - F_x)* (\rho^{2^{-n - 1}} - \rho^{2^{-n}})(z)
\end{equation*}

Just as in the proof of the technical lemma \ref{lem:technique1}, 
we define $\hat{\phi} \coloneq \frac{1}{\int \phi} \sum_{i=0}^{r-1} c_i \phi^{\lambda_i}$ for suitable constants $(c_i),\;(\lambda_i)$, and $\rho \coloneq \hat{\phi}^2*\hat{\phi}$, so that we can write the increment in convolution form~: $\rho^{2^{-m-1}} - \rho^{2^{-m}} = (\rho^{\frac{1}{2}} - \rho)^{2^{-m}} = \hat{\phi}^{2^{-m}}*\check{\phi}^{2^{-m}}$ where $\check{\phi}\coloneq \hat{\phi}^{\frac{1}{2}} - \hat{\phi}^2$.
Recall that $\int \hat{\phi} = 1$ so that $\int \rho = 1$ : we can use it as our mollifier.
Additionaly, $\forall k, 1\leq \abs{k}\leq r-1, \quad \int x^k \hat{\phi}(x) \dd x = 0$ where we fixed $r\in N, \; r>-\alpha$, and the same goes trivially for $\check{\phi}$.
Finally, recall that the coherence bound is verified with $\hat{\phi}$.

\begin{align*}
    \forall z\in \R^d, (f_{x,n+1} - f_{x,n})(z) &= (F_z - F_x)*(\hat{\phi}^{2^{-n}}*\check{\phi}^{2^{-n}})(z)\\
    &= (F_z - F_x)((\hat{\phi}^{2^{-n}}*\check{\phi}^{2^{-n}})(z-\cdot)) \\
    &= (F_z - F_x)*\tilde{\hat{\phi}}^{2^{-n}} ((\check{\phi}^{2^{-n}})(z-\cdot)) \\
    &= \int (F_z - F_x)(\hat{\phi}^{2^{-n}}_y) (\check{\phi}^{2^{-n}})(z-y) \dd y \\
    &= \int (F_y - F_x)(\hat{\phi}^{2^{-n}}_y) (\check{\phi}^{2^{-n}})(z-y) \dd y  \quad (\eqcolon g'_{x,n}(z))\\
    & \quad \quad  + \int (F_z - F_y)(\hat{\phi}^{2^{-n}}_y) (\check{\phi}^{2^{-n}})(z-y) \dd y \quad (\eqcolon g''_{x,n}(z))\\
\end{align*}

We can estimate the first integral with the coherence bound.
Let $\psi \in \D$ be any test-function. Fix a compact $K$ such that $\operatorname{Supp}(\psi)\subset K$.
Then $\operatorname{\check{\phi}^{2^{-n}}*\tilde{\psi}(y)} \subset \bar{K}_1$ where $\bar{K}_1 \coloneq \{x \in \R^d \vert d(x,K)\leq 1\}$ is compact.
We can compute :
\begin{align*}
    \abs{g'_{x,n}(\psi)} & = \abs{\int \dd z \psi(z) \int (F_y - F_x)(\hat{\phi}^{2^{-n}}_y) (\check{\phi}^{2^{-n}})(z-y) \dd y} \\
    & = \abs{\int \dd y (F_y - F_x)(\hat{\phi}^{2^{-n}}_y) (\check{\phi}^{2^{-n}})*\tilde{\psi}(y)} \quad (\text{Fubini, justified a posteriori.}) \\
    & \lesssim \norm{F}_{\alpha, \gamma ; \phi} \int_{\bar{K}_1} \dd y 2^{-n\alpha}(2^{-n}+ \abs{x - y})^{\gamma - \alpha} \abs{\check{\phi}^{2^{-n}}*\tilde{\psi}(y)} \\
    &\lesssim 2^{-n\alpha}\norm{F}_{\alpha, \gamma ; \phi} \norm{\check{\phi}^{2^{-n}}*\tilde{\psi}}_\infty \\
    &\lesssim 2^{-n\alpha} 2^{-nr}\norm{F}_{\alpha, \gamma ; \phi}\norm{\psi}_{C^r} < \infty  \quad \text{by lemma \ref{lem:technique2}} \\
\end{align*}
The finiteness of the left-hand side justifies a posteriori the computation of $g'_{x,n}(\psi)$ which didn't make sense a priori because we didn't prove it was locally integrable.

The second term is easier because $(\check{\phi}^{2^{-n}})(z-y)$ forces $\abs{z-y}\leq 2^-n$ : 
\begin{align*}
    \forall z \in \R^d, \quad \abs{g'_{x,n}(z)} &\leq \norm{\check{\phi}}_{L^1} \sup_{\abs{z-y}\leq 2^-n} \abs{(F_z - F_y)(\hat{\phi}^{2^{-n}}_y)} \\
    &\lesssim 2^{-n\alpha}(2\times2^{-n})^{\gamma - \alpha}\norm{F}_{\alpha, \gamma ; \phi} \\
    &\lesssim 2^{-n \gamma}\norm{F}_{\alpha, \gamma ; \phi} \quad (\text{uniformly in z}).
\end{align*}

In turn, we showed that $f_{x,n+1} - f_{x,n}$ is locally integrable and that for any test-function $\psi\in \D$ and $K$ compact such that $\operatorname{Supp}(\psi)\subset K$ :
\begin{equation*}
    \abs{(f_{x,n+1} - f_{x,n})(\psi)} \lesssim \norm{F}_{\alpha, \gamma ; \phi}\norm{\psi}_{C^r} (2^{-n (\alpha +r)} + 2^{-n \gamma})
\end{equation*}

As a result, because $\gamma, r+\alpha >0$, $f_x = f_{x,l} + \sum_{n=l}^{+\infty}(f_{x,n+1} - f_{x,n})$ is a well-defined distribution of order $r$ (the formula is true for any $l\in \N$).
And obviously $(f_{x,k})_k$ converges to $f_x$.
Recall that $R_n(z) = f_{x,n} (z) + F_x*\rho^{2^{-n}}(z)$. By the mollication lemma, $z \mapsto F_x*\rho^{2^{-n}}(z)$ is a smooth function, converging as a distribution to $F_x$.
So $R_n$ is a locally integrable function converging as a distribution to a limit noted $\mathcal{R}(F)$.
We have the formula :

\begin{equation*}
    \forall l\in \N,\quad  \mathcal{R}(F) = F_x + f_{x,l} + \sum_{n=l}^{+\infty}(f_{x,n+1} - f_{x,n}).
\end{equation*}

We still have to prove the reconstruction bound.
Let $\psi \in B^r$, $\lambda >0$ and $l\in \N$,

\begin{align*}
    f_{x,l}(\psi_x^\lambda) &= \int \dd z f_{x,l}(z) \psi^\lambda(z - x) \\
    &= \int \int \dd z  \dd y (F_z - F_x)*(\hat{\phi}^{2^{-l}})(y)  \hat{\phi}^{2^{-l+1}}(z-y) \psi^\lambda(z - x) \\
\end{align*}

So by the coherence bound, for all $l\in \N$ such that $l\geq \operatorname{log}_2 (\lambda)$,

\begin{align*}
    \abs{f_{x,l}(\psi_x^\lambda)} &\lesssim \norm{F}_{\alpha,\gamma;\phi} \int \int \dd z  \dd y 2^{-l\alpha}(2^{-l}+\abs{x - z})^{\gamma - \alpha}  \abs{\hat{\phi}^{2^{-l+1}}(z-y)} \abs{\psi^\lambda(z - x)} \\
    &\lesssim \norm{F}_{\alpha,\gamma;\phi} \lambda^{\alpha} \int \int \dd z  \dd y (2^{-l}+\abs{x - z})^{\gamma - \alpha}  \abs{\hat{\phi}^{2^{-l+1}}(z-y)} \abs{\psi^\lambda(z - x)} \\
\end{align*}

$\hat{\phi}^{2^{-l+1}}$ is supported in $B(0,2^{-l})$, so that on the domain of the integral, $\abs{x - z}\leq \lambda$. As a result :

\begin{align*}
    \abs{f_{x,l}(\psi_x^\lambda)} &\lesssim \norm{F}_{\alpha,\gamma;\phi} \lambda^{\gamma} \int \int \dd z  \dd y \abs{\hat{\phi}^{2^{-l+1}}(z-y)} \abs{\psi^\lambda(z - x)} \\
    &\lesssim \norm{F}_{\alpha,\gamma;\phi} \lambda^{\gamma} \norm{\hat{\phi}^{2^{-l+1}}}_{L^1} \norm{\psi^\lambda}_{L^1} \\
    &\lesssim \norm{F}_{\alpha,\gamma;\phi} \lambda^{\gamma} \norm{\hat{\phi}}_{L^1} \norm{\psi}_{L^1} \\
    &\lesssim \norm{F}_{\alpha,\gamma;\phi} \lambda^{\gamma} 
\end{align*}

Taking $l\rightarrow +\infty$, we recover the reconstruction bound.










\end{proof}
